% 第三部分 B：摘要（由 main_zh.tex 引入）。
当验证器只能给 GRPO run 的部分样本打标签时，可以由评判模型为其余样本补全奖励。我们沿用第二部分的设置，在 MATH 上用 Qwen3.5-4B-Base 研究这一情形：验证器每组 8 个解答只标注 2 个，另外 6 个或者被丢掉（\texttt{rlvr\_sub}），或者用冻结的 \kev{} 的判断打分（\texttt{sub\_kev}），或者用指针头在迄今收集的标签上在线重新拟合后的 \kev{} 打分（\texttt{sub\_kev\_recal}），或者用一个形状相同、在相同特征和标签上从头训练的输出头打分（\texttt{sub\_probe}）；每个条件两个种子，各 200 步。冻结的评判模型对未标注解答排序良好（AUROC \NUM{fillauc_all_subkev}），但过于慷慨，它的两个策略都学会了在答完之后重复自己，而评判模型会为此付费；它们在 MATH-500 上最终比丢掉这些解答低 \NUM{d_subkev_math} \ci{\NUM{dlo_subkev_math}}{\NUM{dhi_subkev_math}} 个百分点。重新拟合输出头修正了它的校准（期望校准误差 \NUM{fillece_all_subrecal} 对 \NUM{fillece_all_subkev}），但两个策略中仍有一个重复自己。从头训练的输出头在 MATH-500 上最终比丢掉高 \NUM{d_subprobe_math} \ci{\NUM{dlo_subprobe_math}}{\NUM{dhi_subprobe_math}} 个百分点，在 AIME24 上高 \NUM{d_subprobe_aime24}，也高于满标注的基线。它的两个策略都改成了以纯文本的非正式草稿开头；在本项目的其他 GRPO run 中，这一切换无论奖励为何都标志着最好的 run，因此两个种子无法说明增益中有多少属于探针。代码、日志、每个训练解答以及本文的每个数字都在仓库中。
